\section{Related Work}

\subsection{Embodied and LLM-Based Robot Planning}

Recent embodied and language-conditioned robot-planning work studies how
high-level language intent is translated into executable robot behavior in
simulation or on robots, including closed-loop language feedback, embodied
multimodal models, vision-language-action policies, model-based language
grounding, and systems-first integration pipelines such as Inner Monologue
(Huang et al., 2022), PaLM-E (Driess et al., 2023), RT-2 (Brohan et al.,
2023), VoxPoser (Huang et al., 2023), and OK-Robot (Liu et al., 2024). This
includes language-conditioned skill selection, code-generating control, and
program-like robot planning directions exemplified by SayCan (Ahn et al.,
2022), Code as Policies (Liang et al., 2022), and ProgPrompt (Singh et al.,
2022). That literature makes planner quality visible, but it often leaves less
room to isolate what happens at the boundary between planner output and
dispatch. Our paper is narrower: it holds the Isaac Sim setting fixed and
studies how planner-to-executor contract design changes invalid actions,
recovery, success, and planner/tool overhead.

\subsection{Execution Contracts, Runtime Validation, and Robotics Middleware}

A second line of work addresses execution architecture rather than only planner
quality. Relevant threads include ROS 2 middleware and action semantics
(Macenski et al., 2022; Biggs et al., 2019/2020), planning frameworks such as
PlanSys2 (Martin et al., 2021), behavior-tree execution stacks and libraries
such as BehaviorTree.CPP (Colledanchise and Ogren, 2018;
BehaviorTree.CPP documentation), and task-and-motion planning interfaces such
as PDDLStream and the broader TAMP literature (Garrett et al., 2020; Garrett
et al., 2021). The shared intuition is that planner outputs should not
be treated as executable by default; they need an explicit interface to the
execution stack. The present study is closest to that execution-boundary view.
Its runtime layer is intentionally lightweight: R1 validates an emitted
action or tool call against the active contract and allows one repair attempt.
The contribution is therefore evidence about a small validation-and-retry
mechanism in a controlled simulator setting, not a claim of a complete
middleware stack or end-to-end safety case.

\subsection{Controlled Design Studies and Broader Evaluations}

Methodologically, the paper favors a small controlled matrix over benchmark
breadth. Broader embodied-agent evaluations such as EmbodiedBench, IS-Bench,
and the Isaac Sim-based Mind and Motion Aligned benchmark are valuable because
they expose long-horizon execution difficulty, safety failures, or integrated
planning-and-control behavior across larger task sets or evaluation matrices
(Yang et al., 2025; Lu et al., 2025; Kachaev et al., 2025). By varying only
the action interface and runtime validation policy within a fixed Isaac Sim
setting, we instead use a narrower study design to attribute differences in
invalid actions, recovery, success, and planner/tool overhead to explicit
execution-design choices. This complements broader embodied-agent evaluations
rather than competing with them.
